\clearpage\thispagestyle{empty}
\pdfbookmark[0]{Quellen zur Besprechung der Philosophical Transactions}{philosophical1813-sources}
\begin{center}\Large\bfseries Quellen zur Besprechung der Philosophical Transactions\end{center}

Die Übertragung folgt der Besprechung von Carl Friedrich Gauß in den \emph{Werken}, Bd.~VI (Göttingen, 1874), S.~561--563. Der Text im Sammelband entspricht dem Teil der \emph{Göttingischen gelehrten Anzeigen}, Stück~27 vom 14.~Februar 1814, S.~269--272, der mit Ponds Beobachtungen der Sonnenwenden beginnt. Auf den Titelblättern des besprochenen Bandes steht \emph{Year}, in der Besprechung des Sammelbandes \emph{Jear}.

John Ponds erstes gedrucktes Verzeichnis führt 44 Sterne auf (S.~75--76), sein späteres gedrucktes Verzeichnis 84 (S.~280--303). Die ersten beiden Aufsätze zu den Sonnenwenden stehen auf S.~27--30 und 123--125. Der dritte gibt auf S.~304 das Ergebnis für den Sommer 1813 und eine Korrektur des Ergebnisses für den Sommer 1812 an.

Der Nenner bei der Traglast des Drahts ist im Sammelband undeutlich. Die Übertragung gibt die sichtbare Form wieder. Der Göttinger Erstdruck (S.~271) und der Aufsatz von William Hyde Wollaston (S.~117) geben $1\frac{1}{3}$ Gran an. Für Drähte, die sich für astronomische Instrumente eignen, gibt Wollaston außerdem die Dicken $\frac{1}{4000}$ und $\frac{1}{5000}$ Zoll an (S.~115). Im Aufsatz über das Mikrometer gibt er $\frac{1}{25}$ Zoll als Abstand des Lochs vom Mittelpunkt an (S.~119).

Henry Katers Vergleich von Cassegrain- und Gregory-Teleskopen wurde von Sir Joseph Banks vorgelegt (S.~206--212). Das Exemplar der University of Toronto enthält beide Teile der \emph{Philosophical Transactions} von 1813, gedruckt in London von W. Bulmer and Co. und vertrieben von G. and W. Nicol.

KI-Übertragung des deutschen Originals in heutiges Deutsch.

Erstdruck (\href{https://images.p.sub.uni-goettingen.de/iiif/image/gdz:PPN319721507_1814:00000271/full/1650,/0/default.jpg}{269}, \href{https://images.p.sub.uni-goettingen.de/iiif/image/gdz:PPN319721507_1814:00000272/full/1650,/0/default.jpg}{270}, \href{https://images.p.sub.uni-goettingen.de/iiif/image/gdz:PPN319721507_1814:00000273/full/1650,/0/default.jpg}{271}, \href{https://images.p.sub.uni-goettingen.de/iiif/image/gdz:PPN319721507_1814:00000274/full/1650,/0/default.jpg}{272}).

\href{https://archive.org/details/philosophicaltra103royauoft}{Band 103}.
Pond (\href{https://doi.org/10.1098/rstl.1813.0006}{Sommer 1812}, \href{https://doi.org/10.1098/rstl.1813.0020}{Winter 1812}, \href{https://upload.wikimedia.org/wikipedia/commons/e/ed/Philosophical_transactions._v._1-177%2C_1665-66-1886_%28IA_philosophicaltra103royauoft%29.pdf#page=364}{Sommer 1813}, \href{https://doi.org/10.1098/rstl.1813.0011}{Erstes Verzeichnis}, \href{https://doi.org/10.1098/rstl.1813.0035}{Zweites Verzeichnis}).
Wollaston (\href{https://doi.org/10.1098/rstl.1813.0018}{Feine Drähte}, \href{https://doi.org/10.1098/rstl.1813.0019}{Einlinsenmikrometer}).
Kater (\href{https://doi.org/10.1098/rstl.1813.0028}{Teleskopvergleich}).

\clearpage
